\documentclass[12pt,a4paper]{article}

\usepackage[utf8]{inputenc}
\usepackage{geometry}
\usepackage{graphicx}
\usepackage{hyperref}
\usepackage{listings}
\usepackage{xcolor}
\usepackage{amsmath}
\usepackage{amssymb}
\usepackage{fancyhdr}

\geometry{
    left=3cm,
    right=3cm,
    top=3cm,
    bottom=3cm
}

\hypersetup{
    colorlinks=true,
    linkcolor=black,
    urlcolor=blue,
    citecolor=black
}

\lstset{
    basicstyle=\ttfamily\small,
    keywordstyle=\color{blue}\bfseries,
    commentstyle=\color{gray}\itshape,
    stringstyle=\color{red},
    numbers=left,
    numberstyle=\tiny\color{gray},
    frame=single,
    breaklines=true,
    language=C++
}

\pagestyle{fancy}
\fancyhf{}
\fancyhead[L]{\small Tucil 2 IF2211 Strategi Algoritma}
\fancyhead[R]{\small \thepage}
\renewcommand{\headrulewidth}{0.4pt}

% ── title info ──
\title{
    \vspace{-1cm}
    \textbf{Tugas Kecil 2 IF2211 Strategi Algoritma} \\[0.3cm]
    \large Voxelization 3D menggunakan Octree \\
    \large Semester II Tahun 2025/2026
}
\author{
    Made Branenda Jordhy — 13524026 \\
    % tambah anggota di sini
}
\date{\today}

\begin{document}

\maketitle
\thispagestyle{empty}

% signature jika ada
% \begin{center}
%     \includegraphics[width=0.1\textwidth]{media/signature.png}
% \end{center}

\newpage
\tableofcontents
\newpage

% ═══════════════════════════════════════
\section{Pendahuluan}
% ═══════════════════════════════════════

% TODO: isi pendahuluan

% ═══════════════════════════════════════
\section{Dasar Teori}
% ═══════════════════════════════════════

\subsection{Voxelization}
% TODO: jelaskan apa itu voxelization

\subsection{Octree}
% TODO: jelaskan struktur data octree

\subsection{Divide and Conquer}
% TODO: jelaskan D\&C dalam konteks octree

\subsection{SAT (Separating Axis Theorem)}
% TODO: jelaskan SAT untuk triangle-AABB intersection

% ═══════════════════════════════════════
\section{Algoritma Divide and Conquer}
% ═══════════════════════════════════════

% spec: jelaskan langkah-langkah (bukan hanya pseudocode)

\subsection{Langkah-langkah}
% TODO

\subsection{Pseudocode}
% TODO

% ═══════════════════════════════════════
\section{Implementasi}
% ═══════════════════════════════════════

\subsection{Struktur Program}
% TODO: jelaskan modul-modul (parser, octree, exporter, viewer, stats)

\subsection{Alur Program}
% TODO: jelaskan flow dari input sampai output

% ═══════════════════════════════════════
\section{Hasil dan Analisis}
% ═══════════════════════════════════════

% spec: minimal 5 screenshot input dan output

\subsection{Test Case 1}
% TODO: screenshot + tabel statistik

\subsection{Test Case 2}
% TODO

\subsection{Test Case 3}
% TODO

\subsection{Test Case 4}
% TODO

\subsection{Test Case 5}
% TODO

% ═══════════════════════════════════════
\section{Bonus}
% ═══════════════════════════════════════

\subsection{Concurrency}
% TODO: jelaskan implementasi threading

\subsection{Interactive Viewer}
% TODO: jelaskan software renderer, fitur viewer

% ═══════════════════════════════════════
\section{Kesimpulan}
% ═══════════════════════════════════════

% TODO

% ═══════════════════════════════════════
\section{Link Repository}
% ═══════════════════════════════════════

\url{https://github.com/ethj0r/Tucil2_13524026_13524104}

\end{document}
